% =====================================================================
%  Seance 4 -- B. cours 2 : prediction de defaillance (frames 11--16)
% =====================================================================
\section{Cours 2 --- Prédiction de défaillance}

\begin{frame}{Cours 2 --- Prédire la panne : horizon et étiquette}
\begin{columns}[t]
\begin{column}{0.52\textwidth}
\begin{block}{L'horizon de prédiction, choisi AVANT le modèle}
\begin{itemize}
\item \textbf{« Panne dans $h$ heures »} : l'horizon $h$ n'est pas un
détail technique, c'est une \textbf{décision métier}. Trop court : le
sommaire n'a pas le temps d'agir. Trop long : la prédiction devient
bruitée.
\item Notre dossier : $h = 30$ min --- exactement l'étiquette
\py{incident_30min} construite au TP2 à partir de la colonne
\py{alarme} du \textbf{futur uniquement}.
\end{itemize}
\end{block}
\end{column}
\begin{column}{0.44\textwidth}
\begin{ebox}[Le dataset prêt]
\file{nosys\_features.csv} : 34 560 fenêtres de 10 min
(8 serveurs $\times$ 30 jours), 120 positives (0,35 \%) ;
split temporel : entraînement \textbf{avant} le 22/09 (24 192 lignes,
56 positives), test \textbf{après} (10 368 lignes, 64 positives).
\end{ebox}
\begin{kbox}[Le point de bascule]
La séance 3 détectait les \textbf{anomalies} (l'incident est en cours).
Aujourd'hui on le \textbf{précède} --- et tout change, surtout la
validation.
\end{kbox}
\end{column}
\end{columns}
\end{frame}

\begin{frame}[fragile]{Le piège n°1 : la fuite temporelle}
\begin{columns}[t]
\begin{column}{0.52\textwidth}
\begin{lstlisting}
# FAUX : l'alarme courante dans les features
features = ["cpu_moy", "ram_moy", ..., "alarme"]
# or l'etiquette incident_30min CONTIENT
# l'alarme de la fenetre courante :
# label(t) = max(alarme t..t+30min)
\end{lstlisting}
\begin{wbox}[Ce que donne réellement cette fuite (mesuré sur le test)]
\small
\begin{tabular}{@{}lc@{}}
\toprule
 & P(incident $\le$ 30 min)\\
\midrule
alarme courante sonne (40 fen.) & \textbf{100 \%}\\
pas d'alarme courante (10 328) & 0,23 \%\\
\bottomrule
\end{tabular}
\end{wbox}
\end{column}
\begin{column}{0.44\textwidth}
\begin{block}{Pourquoi c'est grave}
\begin{itemize}
\item La variable « prédit » l'étiquette à 100 \% parce qu'elle
\textbf{est} un morceau de l'étiquette : on ne prédit plus, on
\textbf{constate} le présent.
\item En production, cette information arrive \textbf{avec} l'incident,
pas avant : le modèle « magique » devient inutile.
\item Et il reste \textbf{aveugle} aux 24 incidents sans alarme
courante : rappel plafonné à 62,5 \%.
\end{itemize}
\end{block}
\end{column}
\end{columns}
\end{frame}

\begin{frame}{Validation : découpage temporel strict, JAMAIS de shuffle}
\begin{columns}[t]
\begin{column}{0.52\textwidth}
\begin{block}{La règle du dossier}
\begin{center}
\footnotesize entraînement : \textbf{fenêtre $<$ 22/09 00:00} \quad
$\Big|$ \quad test : \textbf{fenêtre $\ge$ 22/09}
\end{center}
\begin{itemize}
\item Le futur ne sert \textbf{jamais} à entraîner : on simule la mise
en production réelle, où l'on prédit \emph{à partir d'aujourd'hui}.
\item La standardisation, les seuils, les encodages : tout se calcule
\textbf{sur le train} puis s'applique au test.
\end{itemize}
\end{block}
\end{column}
\begin{column}{0.44\textwidth}
\begin{ebox}[KFold aléatoire : mesuré sur le dossier]
\small Rappel moyen en KFold aléatoire 5 plis : \textbf{0,60}
$(0{,}54$ à $0{,}63)$ ; découpage temporel strict : \textbf{0,625}.
\end{ebox}
\begin{wbox}[KFold aléatoire : une structure fausse]
\small
Ici les chiffres sont proches --- et c'est le piège : le KFold entraîne
\textbf{sur le futur} et teste sur le passé. La structure de validation
est fausse ; dès que la série dérive (nouvelle charge, évolution du
parc), le KFold promet des performances qui n'existeront pas.
\end{wbox}
\end{column}
\end{columns}
\end{frame}

\begin{frame}{Le modèle du dossier : GradientBoostingClassifier}
\begin{columns}[t]
\begin{column}{0.5\textwidth}
\begin{block}{Réglage et protocole}
\begin{itemize}
\item 12 variables de signaux (niveaux, dispersion, latence p99,
dynamique, contexte cyclique) --- \textbf{sans} l'alarme ;
\item \py{GradientBoostingClassifier} (300 arbres, profondeur 4,
graine 5534) ;
\item split temporel strict au 22/09, seuil 0,5.
\end{itemize}
\end{block}
\begin{solbox}[Matrice de confusion test (mesurée)]
\footnotesize
\begin{tabular}{@{}lcc@{}}
\toprule
 & prédit 0 & prédit 1\\
\midrule
réel 0 & 10 290 & 14\\
réel 1 & 24 & \textbf{40}\\
\bottomrule
\end{tabular}\\[1mm]
rappel \textbf{62,5 \%} \quad précision \textbf{74,1 \%}
\end{solbox}
\end{column}
\begin{column}{0.46\textwidth}
\begin{center}
\includegraphics[width=\textwidth]{f27_rappel_alertes.png}
\end{center}
\begin{kbox}[Lecture]
Avec 5 alertes par jour, on rattrape \textbf{62,5 \%} des pannes à 30
min --- les 10 290 vrais négatifs évitent des alertes inutiles. La
limite à 62,5 \% n'est pas un mauvais réglage : 24 incidents n'ont
\textbf{aucun signal détectable} dans les métriques seules.
\end{kbox}
\end{column}
\end{columns}
\end{frame}

\begin{frame}{Quiz 2 --- prédiction de défaillance (5 min)}
\begin{qbox}[3 questions]
\begin{enumerate}
\item On ajoute l'alarme courante aux variables d'entrée. Sur le jeu de
test, la précision paraît excellente. Le vrai problème :\\
(a) le modèle est trop lent \quad (b) l'alarme \textbf{est} un morceau
de l'étiquette : on constate au lieu de prédire \quad (c) il manque des
données
\item Le bon découpage pour ce dataset est :\\
(a) 80/20 au hasard \quad (b) entraîner avant le 22/09, tester après
\quad (c) KFold aléatoire 5 plis
\item Avec 5 alertes par jour au point d'opération, le rappel mesuré est :\\
(a) 25 \% \quad (b) 62,5 \% \quad (c) 100 \%
\end{enumerate}
\end{qbox}
\vspace{1mm}
\begin{solbox}[Réponses]
\textbf{1 : b} --- fuite temporelle ; l'alarme sonne \emph{avec}
l'incident, jamais avant. \quad \textbf{2 : b} --- le temps ne se
mélange pas. \quad \textbf{3 : b} --- 40 pannes rattrapées sur 64.
\end{solbox}
\end{frame}
